\documentclass[12 pt, twoside]{amsart}

\usepackage{../current_style} % Must be in a child of recipes

\begin{document}

\title{Ginger Snaps}
\maketitle
\subtitle{an hour plus chilling}

\noindent \desc{These are a simple gingersnap recipe. I originally got it from Rose Beranbaum's ``The Baking Bible'' and it hasn't changed since. I've reduced some of the sugar (1 cup + $\frac{1}{3}$ cup of corn syrup down to just a cup and some water) and increased the ginger (3 tsp up to around 4 or 5 tsp). Otherwise it's pretty as is. I found that these took about 20 minutes to cook as opposed to the stated 10-12 minutes, but my oven might have been running cold.}

\recipe{
\begin{ingredients}
\item 1 stick unsalted butter
\item $\frac{1}{3}$ cup water
\item $2 \frac{3}{4}$ cups all purpose flour
\item 1 cup granulated sugar
\item 4 tsp baking powder
\item 2 tsp baking soda
\item $\frac{1}{2}$ tsp salt
\item 4-5 tsp ground ginger
\item 1 egg and the white of another
\end{ingredients}
}{
\begin{directions}
\item Over low heat melt the butter in a saucepan. Add in the water also. When the butter's almost melted remove from the heat and let cool.
\item Meanwhile, ideally in a stand mixer, mix the flour, sugar, baking powder, baking soda, salt, and ginger together.
\item Add the butter mixture to the flour and mix until well combined. Whisk the eggs and add them as well. 
\item Split the dough into three parts and fridge for 30 minutes to a day. When ready to cook again, remove slightly beforehand to allow the dough to warm.
\item Preheat the oven to 350 and tear off chunks of the dough, rolling them into small smooth balls. Set on a baking sheet with oil or parchment paper and bake until cracks start to form and the cookies are golden.
\item When removing the cookies, give them a few minutes to harden and cool before eating.
\end{directions}
}

\end{document}
